\documentclass[10pt,a4paper]{amsart}
\usepackage[margin=19mm]{geometry}
\usepackage{amsmath,amssymb,mathtools}
\usepackage[scr=rsfs,cal=euler]{mathalfa}
\usepackage{xfrac}
\usepackage[unicode,hidelinks,bookmarks=false]{hyperref}
\newcommand{\ie}{i.e.\ }
\newcommand{\cf}{cf.\ }
\newcommand{\resp}{respectively\ }
\newcommand{\Subsection}{\subsection}
\providecommand{\nobreakdash}{\nobreak\hbox{-}}
\setlength{\emergencystretch}{2em}
\multlinegap0pt
\usepackage{amssymb}
\usepackage{mathtools}
\usepackage{bm}
\usepackage[shortlabels]{enumitem}
\usepackage{xcolor}
\newtheorem{theorem}{Theorem}
\newcommand{\bk}{\bm{k}}
\title{\bfseries An Efficient Parity-Blocked Method for Band-Structure Computation of 3D Anisotropic Phononic Crystals}
\author{Jingkai Zhang, Xing-Long Lyu, Tiexiang Li, Wen-Wei Lin}
\date{Source: arXiv:2606.03599v1 (CC BY 4.0)}
\begin{document}
\maketitle

\noindent\emph{Source and licence.} \bfseries An Efficient Parity-Blocked Method for Band-Structure Computation of 3D Anisotropic Phononic Crystals. Version arXiv:2606.03599v1 (CC BY 4.0), distributed by arXiv under the Creative Commons Attribution 4.0 International licence, http://creativecommons.org/licenses/by/4.0/, as recorded in the arXiv licence metadata for this exact version and shown on the abstract page https://arxiv.org/abs/2606.03599v1. Full licence notice: \texttt{licenses/arxiv-2606.03599-CC-BY-4.0.txt}.\par\smallskip

\noindent\textit{Editorial scope.} Counted unit: the article's theorem thm:block\_invariance (source0.tex lines 465--474). The article's complete original proof of it is delivered with it (source0.tex lines 476--511, one \texttt{proof} environment of 36 lines). No mathematics is written, repaired or re-derived: the statement and the proof are the article's own lines, in the article's own notation, and no retained line is substituted or re-flowed.  No further source-local context is needed: every cross-reference of every retained line resolves inside the two delivered spans.  Nothing was omitted from any retained span, so the package carries no omission record.  No retained line contains a citation, so the wrapper carries no bibliography and no bibliography entry.  No retained line contains a figure environment, an image-inclusion command or drawing code, so no image is delivered and the package declares no asset dependency.  Editorial formatting only, in addition: an AMS \texttt{amsart} preamble that restates the article's own declarations and macros used by the retained lines, namely the environments \texttt{theorem} on separate counters, together with the standard packages those lines need.  The retained environments print in this document as Theorem 1 in order of appearance, whereas the article prints the same items with the numbering of its own full document.  The complete native source of the article as distributed in the arXiv e-print for this exact version is bundled unchanged as \texttt{sources/201987/source0.tex}.
\begin{theorem}[Block invariance on even grids]\label{thm:block_invariance}
Assume that \(n_1,n_2,n_3\) are even and that the stiffness matrix \(C_h\) and the mass matrix \(M_h\) are nodewise local multiplication matrices. Then the discrete generalized eigenvalue pair \((A_h(\bk),M_h)\) is compatible with the lifted parity operators:
\begin{equation}\label{eq:A_M_commute_main}
A_h(\bk)\Gamma_a^{(v)}=\Gamma_a^{(v)}A_h(\bk),
\qquad
M_h\Gamma_a^{(v)}=\Gamma_a^{(v)}M_h,
\qquad a\in\{12,13\}.
\end{equation}
Consequently, the velocity space admits an orthogonal decomposition into four common eigenspaces of \((\Gamma_{12}^{(v)},\Gamma_{13}^{(v)})\), and both \(A_h(\bk)\) and \(M_h\) leave each subspace invariant.
\end{theorem}

\begin{proof}
Let \(S_d\) be the Bloch shift with step \(d=(d_1,d_2,d_3)\in\mathbb Z^3\). Since
\[
g_{12}(x+d)=g_{12}(x)(-1)^{d_1+d_2},\qquad
g_{13}(x+d)=g_{13}(x)(-1)^{d_1+d_3},
\]
we have
\begin{equation}\label{eq:Gamma_shift_comm_maintext}
\Gamma_{12}S_d=(-1)^{d_1+d_2}S_d\Gamma_{12},
\qquad
\Gamma_{13}S_d=(-1)^{d_1+d_3}S_d\Gamma_{13}.
\end{equation}
The four body-diagonal shifts used in this paper all have steps \(d=(1,\pm1,\pm1)\), so \(d_1+d_2\) and \(d_1+d_3\) are even. Hence
\[
\Gamma_a S_{d_\ell}=S_{d_\ell}\Gamma_a,
\qquad a\in\{12,13\},\quad \ell=1,2,3,4.
\]
Since \(D_{d_\ell}\) is linear in \(S_{d_\ell}\), and since each reconstructed derivative \(\widetilde D_m\) is a linear combination of \(D_{d_\ell}\), it follows that
\begin{equation}\label{eq:Gamma_Dtilde_comm_main}
\Gamma_a\widetilde D_m=\widetilde D_m\Gamma_a,
\qquad a\in\{12,13\},\quad m=1,2,3.
\end{equation}
The assembly of \(B_h(\bk)\) gives
\begin{equation}\label{eq:B_Gamma_lift_comm_main}
\Gamma_a^{(v)}B_h(\bk)=B_h(\bk)\Gamma_a^{(e)},
\qquad
B_h(\bk)^H\Gamma_a^{(v)}=\Gamma_a^{(e)}B_h(\bk)^H.
\end{equation}
On the other hand, \(C_h\) acts as a local \(6\times6\) Voigt stiffness block at each grid point, and \(M_h\) acts as a local mass block. Neither changes the parity label of a node. Thus
\begin{equation}\label{eq:C_M_Gamma_comm_main}
\Gamma_a^{(e)}C_h=C_h\Gamma_a^{(e)},
\qquad
\Gamma_a^{(v)}M_h=M_h\Gamma_a^{(v)}.
\end{equation}
Combining \eqref{eq:B_Gamma_lift_comm_main} and \eqref{eq:C_M_Gamma_comm_main} yields \eqref{eq:A_M_commute_main}.
\end{proof}

\end{document}
